\documentclass[12pt,a4paper]{article}

% -----------------------------------------------------
% PACKAGES
% -----------------------------------------------------
\usepackage[utf8]{inputenc}
\usepackage[T1]{fontenc}

\usepackage{amsmath,amssymb,amsfonts,mathrsfs}
\usepackage{tikz}
\usetikzlibrary{patterns,arrows,shapes.geometric,fit}

\usepackage{enumitem}
\usepackage{multicol}
\usepackage{lmodern}
\usepackage{times}
\usepackage{pifont}

\usepackage[left=1cm,right=1cm,top=1cm,bottom=1.7cm]{geometry}

\usepackage[colorlinks=true, linkcolor=blue, urlcolor=blue, citecolor=blue]{hyperref}

\usepackage{array,multirow}
\usepackage[most]{tcolorbox}
\tcbuselibrary{skins,hooks}

\usepackage{varwidth}
\usepackage{float}

\usepackage{fancyhdr}
\usepackage{eso-pic}

\usepackage{tkz-tab}
\usepackage{systeme}
\usepackage{verbatim}
\usepackage{color,soul}

% -----------------------------------------------------
% TEXTE EN ARRIÈRE-PLAN (1 seule fois)
% -----------------------------------------------------
\AddToShipoutPicture{
    \AtTextCenter{%
        \makebox[0pt]{\rotatebox{80}{\textcolor[gray]{0.7}{\fontsize{5cm}{5cm}\selectfont PGB}}}
    }
}

% -----------------------------------------------------
% COMMANDES PERSONNALISÉES
% -----------------------------------------------------

% Items
\newcommand\circitem[1]{%
\tikz[baseline=(char.base)]{
\node[circle,draw=gray, fill=red!55,
minimum size=1.2em,inner sep=0] (char) {#1};}}

\newcommand\boxitem[1]{%
\tikz[baseline=(char.base)]{
\node[fill=cyan,
minimum size=1.2em,inner sep=0] (char) {#1};}}

\setlist[enumerate,1]{label=\protect\circitem{\arabic*}}
\setlist[enumerate,2]{label=\protect\boxitem{\alph*}}

\everymath{\displaystyle}

% EXEMPLES
\newcounter{exemple}
\newcommand{\exemple}{%
  \refstepcounter{exemple}%
  \textbf{\textcolor{green}{Exemple \theexemple :}} \ignorespaces
}

% EXERCICES D’APPLICATION
\definecolor{myorange}{rgb}{1.0, 0.8, 0.0}

\newcounter{exerciceapp}
\newcommand{\exerciceapp}{%
  \refstepcounter{exerciceapp}%
  \textbf{\textcolor{myorange}{Exercice d'application \theexerciceapp :}} \ignorespaces
}

% CORRECTIONS
\newcounter{correction}
\newcommand{\correction}{%
  \refstepcounter{correction}%
  \textbf{\textcolor{myorange}{Correction \thecorrection :}} \ignorespaces
}

% Soulignement coloré
\newcommand{\myul}[2][black]{\setulcolor{#1}\ul{#2}\setulcolor{black}}

% Tabulation
\newcommand\tab[1][1cm]{\hspace*{#1}}

% -----------------------------------------------------
% DOCUMENT
% -----------------------------------------------------
\begin{document}

\begin{center}
    \Large\textbf{\underline{Limites et Continuité}}\\[-0.1cm]
    \normalsize\textbf{Prof : M. BA} \hfill \textbf{Classe : 1erL}\\[-0.1cm]
    \textbf{Année scolaire : 2025 -- 2026}
\end{center}

\section*{Introduction aux Polynômes}

\section*{\underline{\textbf{\textcolor{red}{I.  }}}}
\subsection*{\underline{\textbf{\textcolor{red}{1. Monômes}}}} 
\underline{\textbf{\textcolor{red}{Définition et vocabulaire}}}

++++++++++++++++
\section*{I.\quad Limite d’une fonction en un nombre réel}

\subsection*{1)\quad Limite finie d’une fonction en un nombre réel}

\subsubsection*{a)\quad Notion de limite finie en un nombre réel}

Soit $f$ la fonction définie par $f(x)=x^2$. Remplissons le tableau suivant :

\begin{center}
\begin{tabular}{|c|c|c|c|c|c|}
\hline
$x$ & $1{,}999$ & $1{,}9999$ & $1{,}99999$ & $2{,}0001$ & $2{,}00001$ \\
\hline
$f(x)$ & $3{,}996001$ & $3{,}99960001$ & $3{,}9999600001$ & $4{,}00040001$ & $4{,}0000400001$ \\
\hline
\end{tabular}
\end{center}

\bigskip

\begin{center}
\textbf{Exploitation}
\end{center}

D’après le tableau ci-dessus, on constate que si les valeurs de $x$ sont très proches de $2$ alors celles de $f(x)$ sont très proches de $4$.

On dit que la limite de $f(x)$ lorsque $x$ tend vers $2$ est égale à $4$.

On note :
\[
\lim_{x \to 2} f(x)=4
\]

et on lit : « limite de $f(x)$ lorsque $x$ tend vers $2$ égale à $4$. »

\subsubsection*{b)\quad Définition intuitive et notation}

Si les valeurs de $x$ sont très proches d’un nombre réel $a$ (avec $x \neq a$) et que les valeurs correspondantes de $f(x)$ sont très proches d’un nombre réel $l$, alors on dit que la limite de $f(x)$ lorsque $x$ tend vers $a$ est égale à $l$.

On note :
\[
\lim_{x \to a} f(x) = l
\]

et on lit : « limite de $f(x)$ lorsque $x$ tend vers $a$ égale à $l$. »

\bigskip


\subsubsection*{c)\quad Propriété}
\textbf{Propriété :} Lorsqu’une fonction admet une limite en un nombre réel, alors cette limite est unique.

\textbf{Propriété :} 
Si $f$ est une fonction polynôme et si $a$ est un nombre réel, alors
\[
\lim_{x\to a} f(x)=f(a).
\]

\bigskip

\checkmark\ \textbf{Exemple}

Soit $f$ définie par 
\[
f(x)=2x^3-3x^2-x+7.
\]

Calculons la limite de $f(x)$ lorsque $x$ tend vers $-1$.

Comme $f$ est une fonction polynôme, on a :
\[
\lim_{x\to -1} f(x)=f(-1).
\]

Calculons $f(-1)$ :
\[
f(-1)=2(-1)^3-3(-1)^2-(-1)+7
\]
\[
= -2 - 3 + 1 + 7
\]
\[
= 3.
\]

Donc,
\[
\boxed{\lim_{x\to -1} f(x)=3.}
\]

\subsection*{2)\quad Limite infinie d’une fonction en un nombre réel}

\subsubsection*{a)\quad Limite à gauche et limite à droite d’une fonction en un nombre réel}

Soit $f$ la fonction définie par :
\[
f(x)=\frac{1}{x}.
\]

Remplissons le tableau suivant :

\begin{center}
\begin{tabular}{|c|c|c|c|c|c|c|c|c|}
\hline
$x$ & $-0{,}001$ & $-0{,}0001$ & $-0{,}00001$ & $-0{,}000001$ 
& $0{,}001$ & $0{,}0001$ & $0{,}00001$ & $0{,}000001$ \\
\hline
$f(x)$ 
& $-1000$ 
& $-10000$ 
& $-100000$ 
& $-1000000$ 
& $1000$ 
& $10000$ 
& $100000$ 
& $1000000$ \\
\hline
\end{tabular}
\end{center}

\textbf{Interprétation du tableau :}

\medskip

En examinant les quatre premières colonnes du tableau, on constate que lorsque $x$ se rapproche de $0$ tout en restant inférieur à $0$, les valeurs de $f(x)$ deviennent négatives et très grandes en valeur absolue.

On dit alors que :
\[
\lim_{x\to 0^-} f(x) = -\infty.
\]

On lit : « limite de $f(x)$ lorsque $x$ tend vers $0$ par la gauche égale moins l’infini. »

\medskip

En examinant les quatre dernières colonnes du tableau, on constate que lorsque $x$ se rapproche de $0$ tout en restant supérieur à $0$, les valeurs de $f(x)$ deviennent positives et très grandes.

On dit alors que :
\[
\lim_{x\to 0^+} f(x) = +\infty.
\]

On lit : « limite de $f(x)$ lorsque $x$ tend vers $0$ par la droite égale plus l’infini. »

\subsubsection*{b)\quad Définition intuitive et notation}

\checkmark\ 
Si les valeurs de $x$ se rapprochent de $a$ tout en restant inférieures à $a$ et si les valeurs de $f(x)$ deviennent très grandes en valeur absolue, alors :

\[
\lim_{x\to a^-} f(x) = +\infty
\quad \text{ou} \quad
\lim_{x\to a^-} f(x) = -\infty.
\]

\checkmark\ 
Si les valeurs de $x$ se rapprochent de $a$ tout en restant supérieures à $a$ et si les valeurs de $f(x)$ deviennent très grandes en valeur absolue, alors :

\[
\lim_{x\to a^+} f(x) = +\infty
\quad \text{ou} \quad
\lim_{x\to a^+} f(x) = -\infty.
\]

+++++++++++++++++++
\underline{\exemple}

\underline{\textbf{\textcolor{red}{Remarque}}}

\subsection*{\underline{\textbf{\textcolor{red}{2. Polynômes}}}} 
\underline{\textbf{\textcolor{red}{Définition}}}

\underline{\exemple}

\underline{\textbf{\textcolor{red}{Remarque}}}

\section*{\underline{\textbf{\textcolor{red}{II. Trinômes du second degré }}}}

\subsection*{\underline{\textbf{\textcolor{red}{1. Définition et exemple}}}} 

\subsection*{\underline{\textbf{\textcolor{red}{2. Factorisation de $ax^{2}+bx+c$}}}} 

\subsection*{\underline{\textbf{\textcolor{red}{3. Signe de $ax^{2}+bx+c$}}}}

\underline{\exemple}

\underline{\exemple}

\section*{\underline{\textbf{\textcolor{red}{III. Théorème fondamental des polynômes }}}}

\subsection*{\underline{\textbf{\textcolor{red}{1. Racine d’un polynôme}}}}

\underline{\textbf{\textcolor{red}{Définition}}}

\underline{\exemple}

\textbf{\exerciceapp}

\subsection*{\underline{\textbf{\textcolor{red}{2. Théorème fondamental}}}}

\underline{\textbf{\textcolor{red}{Activité}}} 

\underline{\textbf{\textcolor{red}{Exploitation de l’activité}}}

\underline{\textbf{\textcolor{red}{Théorème}}}

\subsection*{\underline{\textbf{\textcolor{red}{3. Division euclidienne}}}}

\underline{\exemple}

\textbf{\exerciceapp}

\subsection*{\underline{\textbf{\textcolor{red}{4. Identification des coefficients}}}}
\underline{\exemple}

\underline{\exemple}

\subsection*{\underline{\textbf{\textcolor{red}{5. Tableau de Horner}}}}
\underline{\exemple}

\underline{\exemple}

\section*{\underline{\textbf{\textcolor{red}{VI. Factorisation pour $ n \geq 3 $ }}}}

\underline{\exemple}

\underline{\exemple}

\end{document}
